\documentclass[12pt,aspectratio=169]{beamer}

\usepackage{fontspec}
\setmainfont{Sylfaen}[
  Path=/Users/surenarzumanyan/Library/Fonts/,
  Extension=.ttf,
  UprightFont=sylfaen,
  AutoFakeBold=2.2
]
\setsansfont{Sylfaen}[
  Path=/Users/surenarzumanyan/Library/Fonts/,
  Extension=.ttf,
  UprightFont=sylfaen,
  AutoFakeBold=2.2
]
\usepackage{polyglossia}
\setmainlanguage{armenian}
\newfontfamily\armenianfont{Sylfaen}[
  Script=Armenian,
  Path=/Users/surenarzumanyan/Library/Fonts/,
  Extension=.ttf,
  UprightFont=sylfaen,
  AutoFakeBold=2.2
]
\newfontfamily\armenianfontsf{Sylfaen}[
  Script=Armenian,
  Path=/Users/surenarzumanyan/Library/Fonts/,
  Extension=.ttf,
  UprightFont=sylfaen,
  AutoFakeBold=2.2
]
\usetheme{Madrid}
\usecolortheme{default}
\setbeamertemplate{navigation symbols}{}
\setbeamertemplate{footline}[frame number]

\title{Թեմատիկ վեբ կայքի ստեղծում առաջատար վեբ ծրագրավորմամբ}
\subtitle{Անձնական նպատակների հետևման վեբ հարթակ}
\author{Սուրեն Արզումանյան}
\institute{Եվրասիա միջազգային համալսարան\\Կառավարման և տեղեկատվական տեխնոլոգիաների ամբիոն}
\date{Երևան, 2026}

\begin{document}

\begin{frame}
  \titlepage
\end{frame}

\begin{frame}{Թեմայի արդիականությունը}
  \begin{itemize}
    \item Անձնական նպատակների կառավարումը դարձել է կրթական և մասնագիտական ինքնակազմակերպման կարևոր գործիք։
    \item Օգտատերերը հաճախ օգտագործում են մի քանի տարբեր հավելված՝ նպատակ, առաջադրանք, սովորություն և վերլուծություն կառավարելու համար։
    \item Առաջատար վեբ ծրագրավորումը հնարավորություն է տալիս այդ գործառույթները միավորել մեկ պաշտպանված և հարմար միջավայրում։
  \end{itemize}
\end{frame}

\begin{frame}{Հետազոտական խնդիրը և նպատակը}
  \textbf{Խնդիր.} Նպատակների կառավարումը մասնատված է տարբեր գործիքների միջև, իսկ ամբողջական լուծումները հաճախ կամ բարդ են, կամ սահմանափակ։

  \vspace{0.5cm}
  \textbf{Նպատակ.} Մշակել առաջատար վեբ ծրագրավորման մոտեցումներով թեմատիկ վեբ կայք՝ անձնական նպատակների սահմանման, հետևման և վերլուծության համար։
\end{frame}

\begin{frame}{Առաջադրված խնդիրները}
  \begin{enumerate}
    \item Գնահատել նպատակադրման և արտադրողականության գիտական մոտեցումները։
    \item Համեմատել գործող նպատակների կառավարման թվային լուծումները։
    \item Մշակել կայքի կառուցվածքը, տվյալների մոդելը և հիմնական գործառույթները։
    \item Գնահատել օգտագործման հարմարավետությունն ու ֆունկցիոնալ ամբողջականությունը։
    \item Ապացուցել, որ ինտեգրված լուծումը նվազեցնում է կառավարման մասնատվածությունը։
  \end{enumerate}
\end{frame}

\begin{frame}{Մեթոդաբանությունը}
  \begin{itemize}
    \item Գրականության որակական վերլուծություն՝ SMART, Goal Setting Theory, OKR, UX և web technology աղբյուրներով։
    \item Գործող լուծումների համեմատություն՝ Todoist, Habitica, Strides և Goals on Track։
    \item Կիրառական մշակում՝ իտերատիվ full-stack մոտեցմամբ։
    \item Ֆունկցիոնալ, responsive և կառուցվածքային գնահատում։
  \end{itemize}
\end{frame}

\begin{frame}{Տեխնոլոգիական լուծումը}
  \begin{itemize}
    \item Next.js 15 App Router և TypeScript՝ full-stack կառուցվածքի համար։
    \item PostgreSQL և Prisma՝ հարաբերական, type-safe տվյալների մոդելի համար։
    \item Auth.js v5՝ Google և email/password նույնականացման համար։
    \item Tailwind CSS, shadcn/ui և Recharts՝ responsive միջերեսի և վերլուծական գրաֆիկների համար։
    \item Docker Compose՝ կրկնելի տեղակայման միջավայրի համար։
  \end{itemize}
\end{frame}

\begin{frame}{Մշակված համակարգի հիմնական գործառույթները}
  \begin{itemize}
    \item Նպատակների ստեղծում, խմբագրում, ջնջում, ֆիլտրում և որոնում։
    \item Ենթանպատակներ, վերջնաժամկետներ, առաջնահերթություններ և կատեգորիաներ։
    \item Առաջընթացի գրանցում և ProgressLog պատմություն։
    \item Dashboard՝ ընդհանուր վիճակագրությամբ և մոտակա վերջնաժամկետներով։
    \item Analytics էջ՝ completion rate, streak և category breakdown ցուցանիշներով։
  \end{itemize}
\end{frame}

\begin{frame}{Հիմնական արդյունքները}
  \begin{itemize}
    \item Ստեղծվել է ամբողջական անձնական նպատակների հետևման թեմատիկ վեբ կայք։
    \item Տվյալների մոդելը կապում է օգտատիրոջը, նպատակները, կատեգորիաները, ենթանպատակները և առաջընթացի պատմությունը։
    \item Համակարգը միավորում է պլանավորումը, գործողության հետևումը և վերլուծությունը մեկ միջավայրում։
    \item Միջերեսը փորձարկվել է desktop և mobile ներկայացումներով։
  \end{itemize}
\end{frame}

\begin{frame}{Եզրակացություններ}
  \begin{itemize}
    \item Գիտական և գործնական աղբյուրները հաստատում են չափելի նպատակների, փուլային բաժանման և feedback-ի կարևորությունը։
    \item Գործող գործիքների համեմատությունը հիմնավորեց ինտեգրված լուծման անհրաժեշտությունը։
    \item Next.js, Prisma, PostgreSQL և Auth.js համադրությունը բավարար է անվտանգ և ընդլայնելի full-stack կայք ստեղծելու համար։
    \item Առաջարկվող լուծումը նվազեցնում է նպատակների կառավարման մասնատվածությունը։
  \end{itemize}
\end{frame}

\begin{frame}{Առաջարկություններ}
  \begin{itemize}
    \item Ավելացնել անհատականացված շաբաթական վերանայման մեխանիզմ։
    \item Կատարել usability testing իրական օգտատերերի՝ ուսանողների և երիտասարդ մասնագետների մասնակցությամբ։
    \item Ավելացնել կարգավորվող push notification-ներ և calendar integration։
    \item Զարգացնել Progressive Web App կամ native mobile տարբերակ։
  \end{itemize}
\end{frame}

\begin{frame}{Սահմանափակումներ և հետագա զարգացում}
  \begin{itemize}
    \item Չկա լայնածավալ իրական օգտատերային փորձարկում։
    \item Չկան թիմային նպատակներ և native mobile հավելված։
    \item Հետագա փուլերում հնարավոր է ավելացնել AI առաջարկություններ, mentor-based նպատակներ և WCAG 2.2 audit։
  \end{itemize}
\end{frame}

\begin{frame}{Շնորհակալություն}
  \centering
  \Large Հարցեր
\end{frame}

\end{document}
